%% kapa: exam=1 round=35 year=2024 period=2 form=A subject=appraisal_law subject_ko=감정평가관계법규 seq=29 number=29
%% source: Q-Net 감정평가사 기출문제 (cst00309) artlSeq=5213229 file=raw/1008/5213229/2교시 문제지 원본.pdf
%% conversion: kordoc 4.24.51 -> tools/gen1.py (auto); review=proofread-sonnet; flags=none
\begin{question}{29}
공간정보의 구축 및 관리 등에 관한 법령상 등록전환을 신청할 수 있는 경우가 아닌 것은?
\begin{choices}
  \item 「건축법」에 따른 건축신고를 한 경우
  \item 도시ㆍ군관리계획선에 따라 토지를 분할하는 경우
  \item 「산지관리법」에 따른 산지일시사용허가를 받은 경우
  \item 지적도에 등록된 토지가 사실상 형질변경되었으나 지목변경을 할 수 없는 경우
  \item 대부분의 토지가 등록전환되어 나머지 토지를 임야도에 계속 존치하는 것이 불합리한 경우
\end{choices}
\end{question}
\answer{4}
